% GENERATED FILE. Edit canonical lessons, then run npm run generate:books.
% canonical-source-hash: fnv1a64:d4c40794ed2536a2
% canonical-lessons: JA-C97-zaisan, JA-C97-tsukaimichi, JA-C97-shudan, JA-C97-houhou, JA-C97-yarikata

\chapter{Property, Use, Means, Method, Way of doing}
\label{ch:ja-property-use-means-method-way-of-doing}
\hlchaptermodality{97}

\begin{chapteropening}
\textbf{By the end of this chapter:} \emph{I can say property, a use, a means, a method and a way of doing.}

Complete the last lesson of chapter 97: I can say property, a use, a means, a method and a way of doing.
\end{chapteropening}

\section[\texorpdfstring{\ja{ざいさん} (zaisan)}{ざいさん (zaisan)}]{\ja{ざいさん} (zaisan) — property, fortune}

\label{lesson:JA-C97-zaisan}



\begin{quote}
Before the new one: say the Japanese for room to spare, then the Japanese for convenience.
\end{quote}

\subsection*{You'll want to know: \ja{ざいさん}}
\textbf{\ja{ざいさん}} — \emph{zaisan} — \textquotedblleft{}property, fortune\textquotedblright{}.

\begin{grammarlens}[title={What you've built}]
One more word for this chapter.
\end{grammarlens}

\subsection*{Your turn}
\begin{itemize}
\raggedright
  \item \emph{Say it:} \emph{zaisan}
  \item \emph{Say it:} \emph{zaisan}, once more
  \item \emph{Say it:} say \emph{zaisan}
  \item \emph{Recall:} say the Japanese for room to spare, then the Japanese for convenience, then say \emph{zaisan} again
\end{itemize}

\begin{tcolorbox}[breakable,colback=teal!4,colframe=teal!35!black,title={Before you move on}]
What does \ja{ざいさん} mean? (\textquotedblleft{}Property, fortune\textquotedblright{}.) Say it once more.
\end{tcolorbox}
\section[\texorpdfstring{\ja{つかいみち} (tsukaimichi)}{つかいみち (tsukaimichi)}]{\ja{つかいみち} (tsukaimichi) — a use, a purpose}

\label{lesson:JA-C97-tsukaimichi}



\begin{quote}
Before the new one: say the Japanese for convenience, then the Japanese for property.
\end{quote}

\subsection*{You'll want to know: \ja{つかいみち}}
\textbf{\ja{つかいみち}} — \emph{tsukaimichi} — \textquotedblleft{}a use, a purpose\textquotedblright{}.

\begin{grammarlens}[title={What you've built}]
One more word for this chapter.
\end{grammarlens}

\subsection*{Your turn}
\begin{itemize}
\raggedright
  \item \emph{Say it:} \emph{tsukaimichi}
  \item \emph{Say it:} \emph{tsukaimichi}, once more
  \item \emph{Say it:} say \emph{tsukaimichi}
  \item \emph{Recall:} say the Japanese for convenience, then the Japanese for property, then say \emph{tsukaimichi} again
\end{itemize}

\begin{tcolorbox}[breakable,colback=teal!4,colframe=teal!35!black,title={Before you move on}]
What does \ja{つかいみち} mean? (\textquotedblleft{}A use, a purpose\textquotedblright{}.) Say it once more.
\end{tcolorbox}
\section[\texorpdfstring{\ja{しゅだん} (shudan)}{しゅだん (shudan)}]{\ja{しゅだん} (shudan) — a means}

\label{lesson:JA-C97-shudan}



\begin{quote}
Before the new one: say the Japanese for property, then the Japanese for a use.
\end{quote}

\subsection*{You'll want to know: \ja{しゅだん}}
\textbf{\ja{しゅだん}} — \emph{shudan} — \textquotedblleft{}a means\textquotedblright{}.

\begin{grammarlens}[title={What you've built}]
One more word for this chapter.
\end{grammarlens}

\subsection*{Your turn}
\begin{itemize}
\raggedright
  \item \emph{Say it:} \emph{shudan}
  \item \emph{Say it:} \emph{shudan}, once more
  \item \emph{Say it:} say \emph{shudan}
  \item \emph{Recall:} say the Japanese for property, then the Japanese for a use, then say \emph{shudan} again
\end{itemize}

\begin{tcolorbox}[breakable,colback=teal!4,colframe=teal!35!black,title={Before you move on}]
What does \ja{しゅだん} mean? (\textquotedblleft{}A means\textquotedblright{}.) Say it once more.
\end{tcolorbox}
\section[\texorpdfstring{\ja{ほうほう} (houhou)}{ほうほう (houhou)}]{\ja{ほうほう} (houhou) — a method}

\label{lesson:JA-C97-houhou}



\begin{quote}
Before the new one: say the Japanese for a use, then the Japanese for a means.
\end{quote}

\subsection*{You'll want to know: \ja{ほうほう}}
\textbf{\ja{ほうほう}} — \emph{houhou} — \textquotedblleft{}a method\textquotedblright{}.

\begin{grammarlens}[title={What you've built}]
One more word for this chapter.
\end{grammarlens}

\subsection*{Your turn}
\begin{itemize}
\raggedright
  \item \emph{Say it:} \emph{houhou}
  \item \emph{Say it:} \emph{houhou}, once more
  \item \emph{Say it:} say \emph{houhou}
  \item \emph{Recall:} say the Japanese for a use, then the Japanese for a means, then say \emph{houhou} again
\end{itemize}

\begin{tcolorbox}[breakable,colback=teal!4,colframe=teal!35!black,title={Before you move on}]
What does \ja{ほうほう} mean? (\textquotedblleft{}A method\textquotedblright{}.) Say it once more.
\end{tcolorbox}
\section[\texorpdfstring{\ja{やりかた} (yarikata)}{やりかた (yarikata)}]{\ja{やりかた} (yarikata) — a way of doing}

\label{lesson:JA-C97-yarikata}



\begin{quote}
Before the new one: say the Japanese for a means, then the Japanese for a method.
\end{quote}

\subsection*{You'll want to know: \ja{やりかた}}
\textbf{\ja{やりかた}} — \emph{yarikata} — \textquotedblleft{}a way of doing\textquotedblright{}.

\begin{grammarlens}[title={What you've built}]
One more word for this chapter.
\end{grammarlens}

\subsection*{Your turn}
\begin{itemize}
\raggedright
  \item \emph{Say it:} \emph{yarikata}
  \item \emph{Say it:} \emph{yarikata}, once more
  \item \emph{Say it:} say \emph{yarikata}
  \item \emph{Recall:} say the Japanese for a means, then the Japanese for a method, then say \emph{yarikata} again
\end{itemize}

\begin{tcolorbox}[breakable,colback=teal!4,colframe=teal!35!black,title={Before you move on}]
What does \ja{やりかた} mean? (\textquotedblleft{}A way of doing\textquotedblright{}.) Say it once more.
\end{tcolorbox}
